\documentclass[10pt,a4paper]{amsart}
\usepackage[margin=19mm]{geometry}
\usepackage{amsmath,amssymb,mathtools}
\usepackage[scr=rsfs,cal=euler]{mathalfa}
\usepackage{xfrac}
\usepackage[unicode,hidelinks,bookmarks=false]{hyperref}
\newcommand{\ie}{i.e.\ }
\newcommand{\cf}{cf.\ }
\newcommand{\resp}{respectively\ }
\newcommand{\Subsection}{\subsection}
\providecommand{\nobreakdash}{\nobreak\hbox{-}}
\setlength{\emergencystretch}{2em}
\multlinegap0pt
\usepackage{amssymb}
\usepackage{mathtools}
\usepackage[shortlabels]{enumitem}
\newtheorem{assumption}{Assumption}
\newtheorem{theorem}{Theorem}
\usepackage{cleveref}
\crefname{assumption}{Assumption}{Assumptions}
\crefname{theorem}{Theorem}{Theorems}
\newcommand{\Dist}{\operatorname{dist}}
\newcommand{\Err}{\mathsf{Err}}
\newcommand{\adm}{\mathrm{adm}}
\newcommand{\bzeta}{\boldsymbol\zeta}
\newcommand{\calA}{\mathcal A}
\newcommand{\calD}{\mathcal D}
\newcommand{\calG}{\mathcal G}
\newcommand{\cl}{\mathrm{cl}}
\newcommand{\comp}{\mathrm{comp}}
\newcommand{\harm}{\mathrm{harm}}
\newcommand{\intg}{\mathrm{int}}
\newcommand{\loc}{\mathrm{loc}}
\newcommand{\proj}{\mathrm{proj}}
\newcommand{\tail}{\mathrm{tail}}
\title{\textbf{Finite-Window Local-to-Clean Transfer and Anti-Phantom Detection\\ for Sharp Navier--Stokes Packages}}
\author{Runlong Yu}
\date{Source: arXiv:2606.18476v1 (CC BY 4.0)}
\begin{document}
\maketitle

\noindent\emph{Source and licence.} \textbf{Finite-Window Local-to-Clean Transfer and Anti-Phantom Detection\\ for Sharp Navier--Stokes Packages}. Version arXiv:2606.18476v1 (CC BY 4.0), distributed by arXiv under the Creative Commons Attribution 4.0 International licence, http://creativecommons.org/licenses/by/4.0/, as recorded in the arXiv licence metadata for this exact version and shown on the abstract page https://arxiv.org/abs/2606.18476v1. Full licence notice: \texttt{licenses/arxiv-2606.18476-CC-BY-4.0.txt}.\par\smallskip

\noindent\textit{Editorial scope.} Counted unit: the article's theorem thm:journal-main (source0.tex lines 2498--2518). The article's complete original proof of it is delivered with it (source0.tex lines 2520--2607, one \texttt{proof} environment of 88 lines). No mathematics is written, repaired or re-derived: the statement and the proof are the article's own lines, in the article's own notation, and no retained line is substituted or re-flowed.  Retained uncounted as the source-local context the proof needs: lines 467--467; lines 471--487; lines 552--565; lines 602--625; lines 2373--2384; lines 2406--2420; lines 2430--2446.  Nothing was omitted from any retained span, so the package carries no omission record.  No retained line contains a citation, so the wrapper carries no bibliography and no bibliography entry.  No retained line contains a figure environment, an image-inclusion command or drawing code, so no image is delivered and the package declares no asset dependency.  Editorial formatting only, in addition: an AMS \texttt{amsart} preamble that restates the article's own declarations and macros used by the retained lines, namely the environments \texttt{assumption}, \texttt{theorem} on separate counters, together with the standard packages those lines need.  The retained environments print in this document as Assumption 1, Theorem 1 in order of appearance, whereas the article prints the same items with the numbering of its own full document.  The complete native source of the article as distributed in the arXiv e-print for this exact version is bundled unchanged as \texttt{sources/203082/source0.tex}.
\section{Structural hypotheses and their status}\label{sec:hypotheses}

\begin{assumption}[Pressure-tail baseline visibility]\label{ass:main-pressure-tail-visibility}
The pressure-natural tail geometry is visible from the older baseline geometry.  More precisely, for the selected finite-window class there are constants and explicit errors such that
\begin{equation}\label{eq:ass-pressure-tail-visibility}
\begin{aligned}
\Dist^{\sharp,\tail}_{\loc,\intg,\tail}
(\calD,\Gamma^{\intg}_{\Lambda,\adm})
&\le
C_{\tail,0}
\Dist_{\loc,\intg,0}
(\calD,\Gamma^{\intg}_{\Lambda,\adm})\\
&\quad
+
\mathfrak E_{\tail,0}(\calD).
\end{aligned}
\end{equation}
Here \(\mathfrak E_{\tail,0}\) consists of projection-tail, harmonic-tail, finite-amplitude, and visibility errors.
\end{assumption}

\begin{assumption}[Dominance of the positive coefficient]\label{ass:main-positive-coefficient}
The clean/chart lower bound dominates the detector loss produced by residual closure and component comparison.  In schematic form,
\begin{equation}\label{eq:ass-positive-coeff}
        c_{\Lambda,0}
        :=
        \mu_\Lambda^{\comp}\lambda_G
        -
        C_{\mathrm{dc}}
        C_{\comp}^{[0,K]}(M_U)
        C_{\comp/0}
        >0.
\end{equation}
If the pressure-tail comparison is inserted into the component-to-baseline comparison, the corresponding pressure-tail constants are included in \(C_{\comp/0}\) and the same positivity condition is used.
\end{assumption}

\begin{theorem}[Pressure-tail closure in baseline gauge]\label{thm:main-pressure-tail-module}
Assume the finite-window baseline visibility, finite-amplitude, same-gauge, projection-tail, and harmonic-tail hypotheses of Appendix A.  Then the pressure-natural tail distance satisfies
\begin{equation}\label{eq:main-pressure-tail-module}
\begin{aligned}
\Dist^{\sharp,\tail}_{\loc,\intg,\tail}
(\calD,\Gamma^{\intg}_{\Lambda,\adm})
&\le
C_{\tail}
\bigl[
(1+C_{\tail/0})
\Dist_{\loc,\intg,0}
(\calD,\Gamma^{\intg}_{\Lambda,\adm})\\
&\qquad
+\delta_0+
\Delta_{\tail/0}
\bigr]
+
\alpha_{\proj}\Delta^{\mathrm{unif}}_{\proj,N}(\calA_\Lambda)
+
\alpha_{\harm}\Delta^{(3/2)}_{\harm,M}.
\end{aligned}
\end{equation}
In particular, after collecting the last three terms into \(\mathfrak E_{\tail,0}(\calD)\), \Cref{ass:main-pressure-tail-visibility} holds.
\end{theorem}

\begin{theorem}[Componentwise residual-ledger closure]\label{thm:main-ledger-module}
Assume the finite-window package geometry, same-chain representative convention, amplitude bound \(M_U\), and near-minimizer condition of Appendix B.  Then
\begin{equation}\label{eq:main-ledger-module}
        \Err_{\comp}^{[0,K]}(\calD;\bzeta_*)
        \le
        C_{\comp}^{[0,K]}(M_U)
        \Dist_{\comp}^{\sharp,[0,K]}(\calD,\calG_{\comp})
        +
        C_{\comp}^{[0,K]}(M_U)
        \delta_{\comp}^{[0,K]}.
\end{equation}
\end{theorem}

\begin{theorem}[Detector comparison]\label{thm:main-detector-module}
Assume the detector-intertwining hypotheses of Appendix C.  Then
\begin{equation}\label{eq:main-detector-module}
\begin{aligned}
        M_{\Lambda}^{\loc}(\calD-\bzeta_*)
        &\ge
        M_{\Lambda}^{\comp}(\Theta_{\Lambda}(\calD-\bzeta_*))
        -
        C_{\mathrm{dc}}
        \Err_{\comp}^{[0,K]}(\calD;\bzeta_*)
        -
        \Delta_{\mathrm{dc}}.
\end{aligned}
\end{equation}
\end{theorem}

\begin{equation}\label{eq:main-detector-after-ledger}
\begin{aligned}
        M_{\Lambda}^{\loc}(\calD-\bzeta_*)
        &\ge
        M_{\Lambda}^{\comp}(\Theta_{\Lambda}(\calD-\bzeta_*))\\
        &\quad
        -
        C_{\mathrm{dc}}C_{\comp}^{[0,K]}(M_U)
        \Dist_{\comp}^{\sharp,[0,K]}(\calD,\calG_{\comp})\\
        &\quad
        -
        C_{\mathrm{dc}}C_{\comp}^{[0,K]}(M_U)
        \delta_{\comp}^{[0,K]}
        -
        \Delta_{\mathrm{dc}}.
\end{aligned}
\end{equation}

\begin{theorem}[Finite-window local-to-clean detection]\label{thm:journal-main}
Fix a finite-window sharp localized Navier--Stokes package \(\calD\), a same-chain representative \(\bzeta_*\), a local-to-clean chart \(\Theta_\Lambda\), and a detector channel \(\Lambda\).  Assume the pressure-tail visibility, componentwise residual-ledger closure, detector comparison, clean gap, chart visibility, component-to-baseline comparison, and positive coefficient hypotheses stated in \Cref{sec:hypotheses}.  Then there are a positive finite-window coefficient \(c_{\Lambda,0}>0\) and an explicit finite-window error functional \(\mathfrak E_{\Lambda,0}(\calD)\) such that
\begin{equation}\label{eq:journal-main}
        M_{\Lambda}^{\loc}(\calD-\bzeta_*)
        \ge
        c_{\Lambda,0}
        \Dist_{\loc,\intg,0}(\calD,\Gamma^{\intg}_{\Lambda,\adm})
        -
        \mathfrak E_{\Lambda,0}(\calD).
\end{equation}
Consequently, if
\begin{equation}\label{eq:positive-detection-threshold}
        \Dist_{\loc,\intg,0}(\calD,\Gamma^{\intg}_{\Lambda,\adm})
        >
        \frac{\mathfrak E_{\Lambda,0}(\calD)}{c_{\Lambda,0}},
\end{equation}
then
\[
        M_{\Lambda}^{\loc}(\calD-\bzeta_*)>0.
\]
\end{theorem}

\begin{proof}
The proof is an assembly of the three modules.

First, \Cref{thm:main-pressure-tail-module} gives pressure-tail closure in baseline gauge.  Whenever the component-to-baseline comparison uses pressure-tail coordinates, this theorem converts those coordinates into the baseline distance and adds only the explicit tail error
\[
        \mathfrak E_{\tail,0}(\calD)
        =
        C_{\tail}(\delta_0+\Delta_{\tail/0})
        +
        \alpha_{\proj}\Delta^{\mathrm{unif}}_{\proj,N}(\calA_\Lambda)
        +
        \alpha_{\harm}\Delta^{(3/2)}_{\harm,M},
\]
up to the finite-window constants already included in \(C_{\comp/0}\).

Second, \Cref{thm:main-ledger-module} gives
\[
        \Err_{\comp}^{[0,K]}(\calD;\bzeta_*)
        \le
        C_{\comp}^{[0,K]}(M_U)
        \Dist_{\comp}^{\sharp,[0,K]}(\calD,\calG_{\comp})
        +
        C_{\comp}^{[0,K]}(M_U)
        \delta_{\comp}^{[0,K]}.
\]
This converts the residual loss in detector comparison into a component distance and a finite-chain near-minimizer error.

Third, \Cref{thm:main-detector-module} gives
\[
\begin{aligned}
        M_{\Lambda}^{\loc}(\calD-\bzeta_*)
        &\ge
        M_{\Lambda}^{\comp}(\Theta_\Lambda(\calD-\bzeta_*))
        -C_{\mathrm{dc}}
        \Err_{\comp}^{[0,K]}(\calD;\bzeta_*)
        -\Delta_{\mathrm{dc}}.
\end{aligned}
\]
Substituting the residual closure yields \eqref{eq:main-detector-after-ledger}.

Fourth, the clean gap and chart visibility give
\[
\begin{aligned}
        M_{\Lambda}^{\comp}(\Theta_\Lambda(\calD-\bzeta_*))
        &\ge
        \mu_\Lambda^{\comp}
        \Dist_{\cl}(\Theta_\Lambda\calD,\Gamma_{\cl,\adm})
        -\Delta_{\cl}\\
        &\ge
        \mu_\Lambda^{\comp}\lambda_G
        \Dist_{\loc,\intg,0}(\calD,\Gamma^{\intg}_{\Lambda,\adm})
        -\Delta_{\cl}
        -\mu_\Lambda^{\comp}\Delta_{\mathrm{chart}}.
\end{aligned}
\]
Fifth, component-to-baseline comparison gives
\[
        \Dist_{\comp}^{\sharp,[0,K]}(\calD,\calG_{\comp})
        \le
        C_{\comp/0}
        \Dist_{\loc,\intg,0}(\calD,\Gamma^{\intg}_{\Lambda,\adm})
        +\Delta_{\comp/0},
\]
with \(C_{\comp/0}\) and \(\Delta_{\comp/0}\) enlarged if the comparison passes through the pressure-tail module.

Collecting the positive and negative coefficients of the baseline distance gives
\[
        c_{\Lambda,0}
        =
        \mu_\Lambda^{\comp}\lambda_G
        -
        C_{\mathrm{dc}}C_{\comp}^{[0,K]}(M_U)C_{\comp/0}.
\]
By \Cref{ass:main-positive-coefficient}, this number is positive.  All remaining terms are placed into
\[
\begin{aligned}
        \mathfrak E_{\Lambda,0}(\calD)
        :=
        &\ \Delta_{\cl}
        +\mu_\Lambda^{\comp}\Delta_{\mathrm{chart}}
        +C_{\mathrm{dc}}C_{\comp}^{[0,K]}(M_U)\Delta_{\comp/0}\\
        &+C_{\mathrm{dc}}C_{\comp}^{[0,K]}(M_U)\delta_{\comp}^{[0,K]}
        +\Delta_{\mathrm{dc}}
        +\mathfrak E_{\tail,0}(\calD),
\end{aligned}
\]
with harmless finite-window constant changes allowed.  This proves \eqref{eq:journal-main}.  The positivity conclusion follows immediately from \eqref{eq:positive-detection-threshold}.
\end{proof}

\end{document}
